\documentclass[12pt, a4paper]{article}

%%% Языковые пакеты %%%
\usepackage[T2A]{fontenc}
\usepackage[utf8]{inputenc}
\usepackage[russian]{babel}

%%% Оформление страницы %%%
\usepackage[left=2.5cm, right=2.5cm, top=2.5cm, bottom=2.5cm]{geometry}
\usepackage{setspace}
\onehalfspacing

%%% Дополнительные пакеты %%%
\usepackage{indentfirst}
\usepackage{graphicx}
\usepackage{hyperref}
\usepackage{csquotes}
\usepackage{enumitem}

\hypersetup{
    colorlinks=true,
    linkcolor=black,
    citecolor=black,
    urlcolor=blue
}

%%% Заголовок %%%
\title{\textbf{Адаптивные топологии мульти-агентных LLM-систем\\с~включением человека в~контур управления}}
\author{}
\date{}

\begin{document}

\maketitle
\thispagestyle{empty}

%%% ===== АННОТАЦИЯ ===== %%%
\begin{center}
\textbf{Аннотация}
\end{center}

Настоящая работа посвящена исследованию влияния коммуникационных топологий мульти-агентных систем на основе больших языковых моделей (LLM) на качество, стоимость и скорость решения сложных задач. Актуальность работы обусловлена стремительным распространением LLM-агентов и отсутствием систематических данных о том, какая структура взаимодействия между ними является оптимальной. Целью исследования является выявление зависимости между типом топологии (звезда, цепочка, полносвязная, дебатная, иерархическая), типом задачи и эффективностью системы, а также разработка адаптивного механизма, способного динамически переключать топологию в~процессе выполнения задачи. Методологически работа строится на контролируемом эксперименте с~варьированием топологии, роли человека и типа задачи. Предполагается, что адаптивное переключение топологий превзойдёт статичные конфигурации, а систематический подход к~позиционированию человека снизит когнитивную нагрузку при сохранении высокого качества решений.

\medskip

\noindent\textbf{Ключевые слова:} мульти-агентные системы, большие языковые модели, топология коммуникации, адаптивные системы, human-in-the-loop, LLM-агенты, оркестрация агентов.

\newpage

%%% ===== ВВЕДЕНИЕ ===== %%%
\section{Введение}

Развитие больших языковых моделей (Large Language Models, LLM) привело к~появлению нового класса интеллектуальных систем~--- мульти-агентных архитектур, в~которых несколько специализированных LLM-агентов кооперируются для решения сложных задач (Wang et al., 2024). Отдельный LLM-агент, как правило, выполняет конкретную роль (исследователь, исполнитель, критик, координатор), тогда как вся система организует их совместную работу посредством определённого протокола коммуникации.

Ключевым, но недостаточно изученным параметром таких систем является \textit{топология коммуникации}~--- граф, описывающий, какие агенты могут обмениваться сообщениями и в~каком порядке. Существующие фреймворки, такие как AutoGen (Wu et al., 2023), CrewAI и ChatDev (Qian et al., 2024), реализуют различные паттерны взаимодействия, однако выбор топологии остаётся, как правило, эвристическим и статичным на протяжении всего процесса решения задачи.

При этом сложные задачи зачастую неоднородны по своей структуре: фаза планирования может потребовать дебатного формата, фаза исполнения~--- чёткой иерархии, а фаза проверки~--- независимой рецензии. Следовательно, фиксированная топология может быть субоптимальной для отдельных фаз, что мотивирует разработку адаптивного механизма переключения.

Ещё одним критически важным фактором является роль человека (human-in-the-loop). В~реальных приложениях человек взаимодействует с~мульти-агентной системой, однако его позиция в~графе коммуникации редко проектируется осознанно. Человек может выступать координатором, рецензентом, судьёй, равноправным участником или наблюдателем, и выбор роли существенно влияет как на качество итогового результата, так и на когнитивную нагрузку.

Таким образом, данное исследование ставит \textbf{цель}: определить, каким образом топология мульти-агентной LLM-системы, механизм её адаптации и позиция человека в~контуре управления влияют на эффективность решения задач различных типов.

Для достижения указанной цели формулируются следующие \textbf{задачи}:
\begin{enumerate}[noitemsep]
    \item Провести сравнительный анализ статичных топологий (звезда, цепочка, полносвязная, дебатная, иерархическая) на задачах разных типов.
    \item Разработать и протестировать механизм динамического переключения топологий между фазами задачи.
    \item Оценить влияние различных ролей и позиций человека на качество решения и когнитивную нагрузку.
    \item Выработать практические рекомендации по проектированию мульти-агентных LLM-систем.
\end{enumerate}

Соответственно, формулируются \textbf{исследовательские вопросы}:
\begin{itemize}[noitemsep]
    \item \textbf{RQ1}: Как различные топологии влияют на качество решений, стоимость (в~токенах) и время выполнения для разных типов задач?
    \item \textbf{RQ2}: Превосходит ли динамическое переключение топологий между фазами задачи статичные конфигурации?
    \item \textbf{RQ3}: Какая роль и позиция человека в~топологии наиболее эффективны?
    \item \textbf{RQ4}: Какая комбинация топологии, адаптации и роли человека обеспечивает оптимальный баланс качества, стоимости и когнитивной нагрузки?
\end{itemize}


%%% ===== ОБЗОР ЛИТЕРАТУРЫ ===== %%%
\section{Обзор литературы}

\subsection{Мульти-агентные системы на основе LLM}

Идея объединения нескольких LLM-агентов для совместного решения задач активно развивается начиная с~2023 года. Wu et al. (2023) представили фреймворк AutoGen, в~котором агенты ведут многоходовые диалоги и могут подключать человека на определённых этапах. Qian et al. (2024) описали ChatDev~--- систему, моделирующую процесс разработки программного обеспечения как взаимодействие ролевых агентов (CEO, программист, тестировщик) в~водопадной модели. Hong et al. (2024) в~проекте MetaGPT предложили стандартизированные операционные процедуры для координации агентов, что снижает количество ошибок за счёт формализации процесса.

Перечисленные работы демонстрируют потенциал мульти-агентных LLM-систем, однако в~каждой из них топология коммуникации задаётся архитектурно, без систематического сравнения альтернативных конфигураций.

\subsection{Топологии коммуникации агентов}

Zhang et al. (2024) предложили G-Designer~--- подход, использующий графовые нейронные сети (GNN) для автоматического проектирования топологии под конкретную задачу. Результаты показали, что оптимальная структура зависит от характера задачи, однако выбор происходит однократно и не пересматривается в~ходе выполнения. Qian et al. (2024b) исследовали закономерности масштабирования (scaling laws) мульти-агентных систем и обнаружили, что топологии типа <<малого мира>> (small-world) демонстрируют повышенную робастность, но изучали только статичные конфигурации.

Liu et al. (2023) в~работе DyLAN реализовали динамический выбор состава команды агентов: на каждом шаге система определяла, какие агенты должны участвовать в~обсуждении. Тем не менее авторы варьировали лишь состав участников, а не саму структуру связей. Zhuge et al. (2024) в~системе GPTSwarm представили оптимизацию графов взаимодействия агентов с~помощью обучения, но без учёта фазовой структуры задач.

\subsection{Безопасность и устойчивость топологий}

Отдельное направление исследований посвящено анализу уязвимостей различных топологий. Yu et al. (2024) в~работе NetSafe продемонстрировали, что определённые конфигурации более подвержены каскадным ошибкам и манипуляциям со стороны скомпрометированного агента. Liang et al. (2024) в~проекте TOMA исследовали многошаговые атаки на мульти-агентные системы с~учётом их топологии. Эти работы важны для контекста данного исследования, поскольку выбор топологии должен учитывать не только производительность, но и надёжность.

\subsection{Роль человека в~мульти-агентных системах}

Включение человека в~контур управления является стандартной практикой в~интерактивных AI-системах (Mosqueira-Rey et al., 2023). Однако в~контексте мульти-агентных LLM-систем систематические исследования позиционирования человека в~графе коммуникации остаются единичными. Wang et al. (2024) в~обзоре LLM-мультиагентных систем выделяют несколько паттернов включения человека, но не предлагают экспериментального сравнения их эффективности.

\subsection{Выводы и исследовательская лакуна}

Анализ литературы выявляет следующий пробел: не существует работы, которая одновременно (a)~систематически сравнивает различные топологии, (b)~изучает механизмы их динамической адаптации в~реальном времени и (c)~экспериментально оценивает влияние позиции человека. Настоящее исследование призвано заполнить эту лакуну.


%%% ===== МЕТОДЫ ===== %%%
\section{Методы исследования}

\subsection{Общий дизайн эксперимента}

Исследование строится как контролируемый эксперимент с~варьированием трёх независимых переменных:
\begin{itemize}[noitemsep]
    \item \textbf{Топология}: звезда (star), цепочка (chain), полносвязная (mesh), дебатная (debate), иерархическая (hierarchical), адаптивная (adaptive).
    \item \textbf{Роль человека}: координатор, рецензент, судья, равноправный участник, наблюдатель.
    \item \textbf{Тип задачи}: программирование, аналитическая задача, креативная задача, задача принятия решений.
\end{itemize}

Для обеспечения сравнимости во всех экспериментах используется фиксированный набор из 3--4 специализированных агентов: \textit{researcher} (поиск и анализ информации), \textit{executor} (генерация решений), \textit{critic} (оценка и критика), \textit{coordinator} (управление процессом).

\subsection{Этапы эксперимента}

Исследование включает четыре последовательных этапа:

\textbf{Этап~1. Baseline.} Тестирование пяти статичных топологий без участия человека. На каждом типе задач каждая топология оценивается по всем метрикам. Цель~--- установить базовую производительность и выявить зависимость <<топология~--- тип задачи>>.

\textbf{Этап~2. Human-in-the-loop.} Для каждой пары <<топология~--- тип задачи>> из Этапа~1 добавляется человек в~каждой из пяти ролей. Цель~--- количественно оценить вклад человека и определить оптимальную роль для каждой конфигурации.

\textbf{Этап~3. Adaptive.} Разработка и тестирование адаптивного механизма, который переключает топологию между фазами задачи (планирование, исполнение, проверка). Роль человека также может изменяться от фазы к~фазе. Цель~--- проверить гипотезу о~превосходстве адаптивного подхода над статичным.

\textbf{Этап~4. Анализ.} Статистическая обработка результатов с~использованием дисперсионного анализа (ANOVA), попарных сравнений (post-hoc тесты) и визуализации. Формулировка практических рекомендаций.

\subsection{Метрики}

Для комплексной оценки используются следующие метрики:
\begin{itemize}[noitemsep]
    \item \textbf{Качество}: точность и полнота решения, оцениваемые комбинированным методом (LLM-as-judge и экспертная оценка человеком).
    \item \textbf{Эффективность}: количество использованных токенов и стоимость API-вызовов.
    \item \textbf{Время}: общая продолжительность выполнения задачи и количество раундов коммуникации.
    \item \textbf{Human factors}: когнитивная нагрузка (NASA-TLX), субъективная удовлетворённость, объём затраченного человеком времени.
\end{itemize}

\subsection{Технический стек}

Реализация системы осуществляется на языке Python~3.10+. В~качестве LLM используются модели OpenAI (GPT-4o, GPT-4o-mini) и Anthropic (Claude 3.5 Sonnet, Claude 3 Haiku). Графы коммуникации между агентами строятся с~помощью фреймворка LangGraph, обеспечивающего декларативное описание потоков взаимодействия. Взаимодействие с~системой осуществляется через интерфейс командной строки (CLI). Для анализа и визуализации результатов используются Pandas и Matplotlib.

Архитектура системы включает следующие компоненты: Agent Layer (набор специализированных агентов), Topology Manager (управление графом связей), Phase Manager (определение и переключение фаз), Orchestrator (центральный координатор), Human Adapter (CLI-интерфейс взаимодействия с~человеком) и Evaluation Framework (сбор метрик).


%%% ===== ПРЕДПОЛАГАЕМЫЕ РЕЗУЛЬТАТЫ ===== %%%
\section{Предполагаемые результаты}

На основании анализа литературы и пилотных наблюдений ожидаются следующие результаты.

\textbf{По~RQ1.} Предполагается, что различные топологии покажут статистически значимые различия в~производительности, причём оптимальная топология будет зависеть от типа задачи. Для задач программирования наиболее эффективной должна стать иерархическая топология (чёткое разделение ролей), для аналитических~--- дебатная (множественные точки зрения), для креативных~--- полносвязная (свободный обмен идеями), для задач принятия решений~--- звёздная (централизованная агрегация).

\textbf{По~RQ2.} Адаптивное переключение топологий между фазами задачи должно превзойти лучшую статичную конфигурацию за счёт возможности использовать сильные стороны каждой топологии на соответствующей фазе.

\textbf{По~RQ3.} Роль человека-рецензента предположительно обеспечит наилучший баланс качества и когнитивной нагрузки для большинства типов задач, тогда как роль координатора будет давать максимальное качество при повышенной нагрузке.

\textbf{По~RQ4.} Оптимальной комбинацией станет адаптивная топология с~человеком в~роли рецензента или наблюдателя, вмешивающегося только на критических фазах.

Практическими результатами работы станут: набор бенчмарков с~результатами на различных типах задач и руководство (design guidelines) для разработчиков по выбору топологии и интеграции человека в~зависимости от типа задачи.


%%% ===== ЗАКЛЮЧЕНИЕ ===== %%%
\section{Заключение}

Данное исследование направлено на решение актуальной научной проблемы~--- систематического изучения коммуникационных топологий мульти-агентных LLM-систем с~учётом динамической адаптации и роли человека. Научная новизна состоит в~том, что впервые будет проведено комплексное сравнение топологий в~контролируемых условиях с~едиными метриками, изучен механизм их адаптации в~реальном времени и экспериментально оценено влияние позиционирования человека в~графе агентов.

Результаты исследования имеют значительный практический потенциал: они могут быть использованы при разработке AI-assisted инструментов для программирования, аналитики и принятия решений, при проектировании collaborative AI систем в~индустрии, а также при оптимизации человеко-машинных команд в~различных предметных областях.


%%% ===== БИБЛИОГРАФИЯ ===== %%%
\newpage
\begin{thebibliography}{99}

\bibitem{qian2024scaling}
Qian, C., Ye, Z., Liu, W., Chen, N., \& Liu, Z. (2024). Scaling large-language-model-based multi-agent collaboration. \textit{arXiv preprint arXiv:2406.07155}.

\bibitem{yu2024}
Yu, M., Wang, S., Zhang, G., \& Li, H. (2024). NetSafe: Exploring the topological safety of multi-agent networks against adversarial attacks. \textit{arXiv preprint arXiv:2410.15686}.

\bibitem{hong2024}
Hong, S., Zhuge, M., Chen, J., Zheng, X., Cheng, Y., Zhang, C., \& Wu, Y. (2024). MetaGPT: Meta programming for a multi-agent collaborative framework. \textit{Proceedings of ICLR 2024}.

\bibitem{liu2023}
Liu, Z., Zhang, Y., Li, P., Liu, Y., \& Yang, D. (2023). Dynamic LLM-agent network: An LLM-agent collaboration framework with agent team optimization. \textit{arXiv preprint arXiv:2310.02170}.

\bibitem{mosqueira2023}
Mosqueira-Rey, E., Hern\'{a}ndez-Pereira, E., Alonso-R\'{i}os, D., Bobes-Bascar\'{a}n, J., \& Fern\'{a}ndez-Leal, \'{A}. (2023). Human-in-the-loop machine learning: A state of the art. \textit{Artificial Intelligence Review}, 56(4), 3005--3054.

\bibitem{qian2024}
Qian, C., Liu, W., Liu, H., Chen, N., Dang, Y., Li, J., \& Liu, Z. (2024). ChatDev: Communicative agents for software development. \textit{Proceedings of ACL 2024}, 15174--15186.

\bibitem{wang2024}
Wang, L., Ma, C., Feng, X., Zhang, Z., Yang, H., Zhang, J., \& Wen, J.-R. (2024). A survey on large language model based autonomous agents. \textit{Frontiers of Computer Science}, 18(6), 186345.

\bibitem{wu2023}
Wu, Q., Bansal, G., Zhang, J., Wu, Y., Li, B., Zhu, E., \& Wang, C. (2023). AutoGen: Enabling next-gen LLM applications via multi-agent conversation. \textit{arXiv preprint arXiv:2308.08155}.

\bibitem{liang2024}
Liang, R., Ye, J., Chen, Z., \& Zhu, Q. (2024). Tipping the dominos: Topology-aware multi-hop attacks on LLM-based multi-agent systems. \textit{arXiv preprint arXiv:2512.04129}.

\bibitem{zhang2024}
Zhang, Y., Xu, K., Li, Y., Liu, Z., \& Shen, D. (2024). G-Designer: Architecting multi-agent communication topologies via graph neural networks. \textit{arXiv preprint arXiv:2410.11782}.

\bibitem{zhuge2024}
Zhuge, M., Wang, W., Kirsch, L., Faccio, F., Khizbullin, D., \& Schmidhuber, J. (2024). GPTSwarm: Language agents as optimizable graphs. \textit{Proceedings of ICML 2024}.

\end{thebibliography}

\bigskip
\noindent\textbf{Word Count:} 2180

\end{document}
